\section{Anforderungen}
Das zu entwickelnde Framework soll eine Python-Schnittstelle enthalten,
da Python eine moderne, weit verbreitete und vor allem gut lesbare Programmiersprache ist, die vielen Studierenden vertraut ist.
Die gute Lesbarkeit und klare Syntax sind für ein interaktives Framework wie dieses von Vorteil \doublecite[35-36]{python3_das_umfassende_handbuch}. 
Das Framework soll für die Nutzung innerhalb der JupyterLab-Umgebung konzipiert sein, da dadurch eine gute Portabilität gewährleistet wird.
Außerdem ist JupyterLab bereits in den Robotik-Übungen an der Fachhochschule Dortmund im Einsatz.

Es wird vorausgesetzt, dass Transformationen von Körpern im dreidimensionalen Raum visualisiert werden können.
Diese Visualisierungen sollen bevorzugt direkt innerhalb von JupyterLab erfolgen.
Sollte das nicht möglich sein, wäre es möglich die visuelle Darstellung in einem separaten Browserfenster vorzunehmen.
Für die Visualisierung soll keine lokale Programminstallation erforderlich sein, da dies die Portabilität einschränken würde.
Darüber hinaus soll das Framework so gestaltet sein, dass es die komfortable Erstellung und Anwendung von Übungsaufgaben unterstützt und vereinfacht.
Zusätzlich wird gefordert, dass das Framework mit \textit{SymPy}\footnote{Sympy ist eine freie Python-Bibliothek für symbolische Mathematik, die es ermöglicht, mathematische Ausdrücke zu vereinfachen oder Gleichungen symbolisch zu lösen.} kompatibel ist, da das Framework in bereits vielen schon existierende
Notebooks mit Übungsaufgaben integriert werden könnte.


Das Framework soll grundlegende Konzepte der räumlichen Transformation abdecken.
Dazu zählen insbesondere Rotationen, Eulerwinkel, Roll-Pitch-Yaw-Winkel sowie Quaternionen zur Darstellung und Analyse von Orientierungen im dreidimensionalen Raum.
Darüber hinaus ist vorgesehen, auch einfache Modelle radgetriebener Roboter zu integrieren, um kinematische Zusammenhänge und Bewegungsabläufe anschaulich darzustellen.

Langfristig ist geplant, das Framework um Funktionalitäten zur Modellierung stationärer Roboter, insbesondere Knickarmroboter, zu erweitern.
In diesem Zusammenhang sollen auch die Vorwärts- und Rückwärtstransformationen behandelt werden.
Aufgrund des Umfangs und der Komplexität dieser Themen wird dieser Teil voraussichtlich im Rahmen einer späteren Bachelorarbeit realisiert und vertieft.


Zwar existieren bereits etablierte Werkzeuge wie beispielsweise Peter Corke’s Robotics Toolbox for Python, welche eine Vielzahl an Funktionen für die Modellierung und Analyse robotischer Systeme bieten.
Allerdings ist die Nutzung dieser Toolbox im Kontext von JupyterLab nur eingeschränkt möglich.
Die dort verwendeten Visualisierungen basieren auf einem separaten \textit{Renderer}\footnote{Ein Renderer ist ein Software- oder Hardware-System,
das zur Erzeugung von Bildern oder Szenen aus einem 3D-Modell verantwortlich ist. Es berechnet, wie Licht mit den Objekten interagiert und wie diese auf der 2D-Oberfläche eines Displays dargestellt werden.
In der Computergrafik werden verschiedene Rendering-Techniken verwendet, um realistische oder stilisierte Darstellungen zu erzeugen.} namens Swift,
welcher ein eigenes Browserfenster zur Darstellung benötigt und somit nicht direkt in JupyterLab eingebettet ist \doublecite[77]{Robotics_Vision_and_Control}.
Darüber hinaus ist die Toolbox sehr umfangreich und adressiert viele Konzepte, die für die Anforderungen und Inhalte dieses Projekts nicht von unmittelbarer Relevanz sind.

Auch Tools wie CoppeliaSim bieten umfangreiche Simulationsmöglichkeiten für Robotersysteme,
inklusive physikbasierter Modellierung, Sensorik und komplexer Steuerungsarchitekturen.
CoppeliaSim ist ein eigenständiger 3D-Robotersimulator mit einer integrierten Skriptumgebung.
Es kann sowohl als eigenständige Anwendung genutzt als auch in eine Hauptanwendung eingebettet werden.
Die Software bietet eine umfangreiche API und eine kompakte Struktur, die eine einfache Integration ermöglicht.
Für die Nutzung ist jedoch eine lokale Installation erforderlich, da CoppeliaSim als Desktop-Anwendung konzipiert ist.
Die Installation umfasst die Einrichtung der Simulationsumgebung sowie der erforderlichen Komponenten für die Skriptverarbeitung und Visualisierung. 
Im Kontext dieses Projekts, das auf eine browserbasierte und plattformunabhängige Lösung abzielt, stellt die Notwendigkeit einer lokalen Installation eine Einschränkung dar.
Zudem ist CoppeliaSim auf die Simulation vollständiger Roboterumgebungen ausgerichtet, was über die Anforderungen dieses Projekts hinausgeht \doublecite{coppeliaSim_manual}.

Eine maßgeschneiderte Lösung, die gezielt auf die Bedürfnisse und Themen des Lehrkonzepts an der Fachhochschule Dortmund zugeschnitten ist, bietet daher klare Vorteile.
Sie erleichtert die Erstellung spezifischer Übungsaufgaben durch den Aufgabensteller,
fördert die Anschaulichkeit durch eingebettete Visualisierungen für Lernende und trägt insgesamt zur Effizienz und Qualität der Lehre im Bereich der robotischen Transformationen bei.

